% eksperymentalne tłumaczenie na włoski

%

\section{Elementarne wyniki} % lang-pl
\section{Risultati elementari} % lang-it
\begin{definition}[asse del segmento] % lang-it
    \label{def_symetralna} % lang-it
    La retta perpendicolare a un segmento e passante per il suo punto medio si chiama asse del segmento. % lang-it
    \index{asse del segmento} % lang-it
\end{definition} % lang-it


\begin{definition}[symetralna] % lang-pl
    \label{def_symetralna} % lang-pl
	Prostą prostopadłą do odcinka i przechodzącą przez jego środek nazywamy symetralną. % lang-pl
    \index{symetralna} % lang-pl
\end{definition} % lang-pl
Symetralna odcinka $AB$ jest miejscem geometrycznym punktów, które są równodległe od obu końców odcinka, a zatem są środkiem pewnego okręgu przechodzącego przez punkty $A$ oraz $B$ (Guzicki \cite[s. 14, 15]{guzicki_2021}). % lang-pl
L'asse del segmento $AB$ è il luogo geometrico dei punti equidistanti dai due estremi del segmento e quindi costituisce l'insieme dei centri delle circonferenze passanti per i punti $A$ e $B$ (Guzicki \cite[s. 14, 15]{guzicki_2021}). % lang-it


W szczególności, trzy współliniowe punkty nie leżą na jednym okręgu, ale już trzy niewspółliniowe punkty wyznaczają taki okrąg jednoznacznie. % lang-pl
In particolare, tre punti allineati non appartengono a una stessa circonferenza, mentre tre punti non allineati determinano univocamente una tale circonferenza. % lang-it


\index{współliniowy} % lang-pl
\index{allineato} % lang-it


Okrąg przechodzący przez wszystkie wierzchołki wielokąta nazywamy opisanym na nim, zatem poprzednie stwierdzenie można wysłowić krótko ,,na każdym trójkącie można opisać dokładnie jeden okrąg''. % lang-pl
La circonferenza passante per tutti i vertici di un poligono si dice circoscritta al poligono; pertanto l'affermazione precedente può essere formulata brevemente come ``per ogni triangolo esiste una e una sola circonferenza circoscritta''. % lang-it


Mówimy też, że wielokąt jest wpisany w okrąg. % lang-pl
Si dice anche che il poligono è inscritto nella circonferenza. % lang-it


\index{okrąg!opisany} % lang-pl
\index{circonferenza!circoscritta} % lang-it

% TODO: https://en.wikipedia.org/wiki/Bisection#Line_segment_bisector
% TODO: https://en.wikipedia.org/wiki/Circumcircle

\begin{proposition}
	\label{hartshorne_52x}
    Niech $AB$ będzie odcinkiem. % lang-pl
    
    Sia $AB$ un segmento. % lang-it

	Istnieje wtedy trójkąt równoramienny, którego podstawą jest $AB$. % lang-pl
    Allora esiste un triangolo isoscele avente $AB$ come base. % lang-it
    
    \index{trójkąt!równoramienny} % lang-pl
    \index{triangolo!isoscele} % lang-it
\end{proposition}

Powyższe stwierdzenie jest ciekawe, bo jest prawdziwe na płaszczyźnie Hilberta: jego prawdziwość nie zależy od aksjomatu Pascha. % lang-pl
L'affermazione precedente è interessante perché è valida nel piano di Hilbert: la sua validità non dipende dall'assioma di Pasch. % lang-it


\index{płaszczyzna!Hilberta} % lang-pl
\index{piano!di Hilbert} % lang-it


\index{aksjomat!Pascha} % lang-pl
\index{assioma!di Pasch} % lang-it


W geometrii nieeuklidesowej może nie istnieć trójkąt równoboczny o danej podstawie. % lang-pl
Nella geometria non euclidea può non esistere un triangolo equilatero con una base assegnata. % lang-it


Ale na płaszczyźnie Hilberta z aksjomatem (P), bez zakładania aksjomatu (E), istnieje trójkąt równoboczny o danej podstawie -- patrz ćwiczenie u Hartshorne'a \cite[s. 184]{hartshorne2000}. % lang-pl
Tuttavia, nel piano di Hilbert con l'assioma (P), senza assumere l'assioma (E), esiste un triangolo equilatero con una base assegnata; si veda l'esercizio in Hartshorne \cite[s. 184]{hartshorne2000}. % lang-it

\begin{proposition}
\label{hartshorne_52}%
	Linia środkowa (odcinek łączący środki pewnych dwóch boków trójkąta) jest równoległa do trzeciego boku. % lang-pl
    Il segmento che congiunge i punti medi di due lati di un triangolo è parallelo al terzo lato. % lang-it
    
    \index{segmento dei punti medi} % lang-it

    \index{linia środkowa} % lang-pl
    
    Jej długość jest dwukrotnie mniejsza od długości tego boku. % lang-pl
    La sua lunghezza è pari alla metà della lunghezza di tale lato. % lang-it
\end{proposition}
% Hartshorne s. 52

Tego stwierdzenia nie ma w Elementach Euklidesa, ale można wyprowadzić je z księgi I (I.29, I.26, I.34), jak wspomina Hartshorne \cite[s. 45, 52. 53]{hartshorne2000}. % lang-pl
Questa affermazione non compare negli \emph{Elementi} di Euclide, ma può essere dedotta dal Libro I (I.29, I.26, I.34), come osserva Hartshorne \cite[s. 45, 52--53]{hartshorne2000}. % lang-it


Patrz też Bogdańska, Neugebauer \cite[s. 15, 52]{neugebauer_2018}. % lang-pl
Si vedano anche Bogdańska e Neugebauer \cite[s. 15, 52]{neugebauer_2018}. % lang-it


Joanna Jaszuńska napisze ,,Niby nic'' w $\Delta_{17}^5$. % lang-pl
Joanna Jaszuńska pubblicherà l'articolo ``Niby nic'' in $\Delta_{17}^5$. % lang-it

% \todofoot{https://www.deltami.edu.pl/2017/05/niby-nic/}

\begin{corollary}
	Niech $ABC$ będzie trójkątem, zaś punkty $D$, $E$ i $F$ środkami jego boków. % lang-pl
    Sia $ABC$ un triangolo e siano $D$, $E$ e $F$ i punti medi dei suoi lati. % lang-it

	Wtedy cztery małe trójkąty utworzone na bokach $DE$, $EF$, $FD$ są przystające do siebie. % lang-pl
    Allora i quattro piccoli triangoli determinati dai segmenti $DE$, $EF$ e $FD$ sono congruenti tra loro. % lang-it
\end{corollary}

Do tego wniosku potrzeba dodatkowo cechy przystawania bok-bok-bok (I.8). % lang-pl
Per ottenere questa conclusione è inoltre necessario utilizzare il criterio di congruenza lato-lato-lato (I.8). % lang-it


\index{cecha przystawania!bok-bok-bok} % lang-pl
\index{criterio di congruenza!lato-lato-lato} % lang-it


Pompe \cite[s. 40]{pompe_2022} uogólnia fakt \ref{hartshorne_52} do czworokątów: % lang-pl
Pompe \cite[s. 40]{pompe_2022} generalizza il fatto \ref{hartshorne_52} ai quadrilateri: % lang-it

\begin{proposition}
	Dany jest czworokąt wypukły $ABCD$ oraz punkty $K$ i $L$ leżące na środkach boków $AD$ i $BC$. % lang-pl
    Siano dati un quadrilatero convesso $ABCD$ e i punti $K$ e $L$, punti medi rispettivamente dei lati $AD$ e $BC$. % lang-it

	Wtedy $KL \le \frac 12 (AB +CD)$ z równością wtedy i tylko wtedy, gdy proste $AB$ i $CD$ są równoległe. % lang-pl
    Allora $KL \le \frac 12 (AB + CD)$, con uguaglianza se e solo se le rette $AB$ e $CD$ sono parallele. % lang-it

	(Wtedy prosta $KL$ też jest równoległa do nich). % lang-pl
    (In tal caso anche la retta $KL$ è parallela ad esse). % lang-it
\end{proposition}

\begin{corollary}
	Prosta przechodząca przez środek ramienia trapezu i równoległa do jego podstaw przechodzi przez środek drugiego ramienia oraz środki obu przekątnych. % lang-pl
    La retta passante per il punto medio di un lato obliquo di un trapezio e parallela alle sue basi passa per il punto medio dell'altro lato obliquo e per i punti medi di entrambe le diagonali. % lang-it
\end{corollary}

\begin{definition}[dwusieczna] % lang-pl
	Niech $\angle ABC$ będzie kątem. % lang-pl
\begin{definition}[bisettrice] % lang-it
    Sia $\angle ABC$ un angolo. % lang-it

	Półprostą $BD$ leżącą w jego wnętrzu taką, że kąty $\angle ABD = \angle DBC$ mają równe miary, nazywamy dwusieczną kąta $\angle ABC$. % lang-pl
    La semiretta $BD$, contenuta nell'interno dell'angolo e tale che gli angoli $\angle ABD$ e $\angle DBC$ abbiano la stessa ampiezza, si chiama bisettrice dell'angolo $\angle ABC$. % lang-it
\end{definition}

Dwusieczna kąta (wypukłego) jest miejscem geometrycznym punktów, które leżą w jego wnętrzu i~są równoodległe od ramion. % lang-pl
La bisettrice di un angolo (convesso) è il luogo geometrico dei punti che giacciono nel suo interno e sono equidistanti dai suoi lati. % lang-it


Zadania o niej znajdą się w $\Delta_{16}^{4}$. % lang-pl
Esercizi su questo argomento compariranno in $\Delta_{16}^{4}$. % lang-it


Symetralne i dwusieczne jako bliźnięta jedno- bądź dwujajowe będą tematem artykułu w $\Delta_{13}^8$. % lang-pl
Gli assi dei segmenti e le bisettrici, considerate come gemelli monozigoti o dizigoti, saranno il tema di un articolo in $\Delta_{13}^8$. % lang-it

% \todofoot{https://www.deltami.edu.pl/2016/04/dwusieczne/}
% \todofoot{https://www.deltami.edu.pl/2013/08/symetralna-i-dwusieczna-bliznieta-jedno-czy-dwujajowe} %  + www.deltami.edu.pl/media/issues/2013/08/delta-2013-08.pdf

\begin{proposition}
	Dwusieczna $t$ kąta przy podstawie trójkąta równoramiennego o podstawie $a$ oraz ramionach $b, b$ spełnia podwójną nierówność % lang-pl
	
	La bisettrice $t$ di un angolo alla base di un triangolo isoscele di base $a$ e lati congruenti $b, b$ soddisfa la doppia disuguaglianza % lang-it
	\begin{equation}
		\sqrt{2} \cdot \frac{ab}{a+b} < t < 2 \cdot \frac{ab}{a+b}.
	\end{equation}
\end{proposition}
% https://en.wikipedia.org/wiki/List_of_triangle_inequalities#Isosceles_triangle

%